%paper also and it is very nice !
\documentclass[a4paper,11pt]{paper}

\usepackage{graphicx}
%\usepackage{float}% If comment this, figure moves
\usepackage{mwe}    % loads »blindtext« and »graphicx«
\usepackage{subfig} % error
\usepackage[english]{babel}
\usepackage{ucs}               % require separate debian package 
\usepackage[utf8x]{inputenc}   % require ucs. Allows both Italian and English letters.
\usepackage{textcomp}          % allow euro symbol
\usepackage{fancyhdr}
\usepackage{amssymb,amsmath}
\usepackage{threeparttable}
\usepackage{color}
\usepackage{url}
\usepackage{natbib}  % substituted hardvard with natbib as the former package doesn't like the hyperref one
\usepackage{rotating} %to have text vertical alignment in some table headers
\usepackage{url}
\usepackage{moreverb}
\usepackage{tabularx} % allow line breaks in tables
\usepackage{html,htmllist} % don't think will be needed
\usepackage{hyperref}

\bibliographystyle{agsmurl}  % bibliographic style


%opening
\title{FFSM++ complete output}

\author{Ahmed Barkaoui, Sylvain Caurla, Philippe Delacote, Antonello Lobianco}

\date{
March 1, 2014
}

\bibliographystyle{agsmurl}

\begin{document}
\maketitle

\begin{abstract}
This book contains the complete results of FFSM++ in terms of tables and charts, hierarchically organised by region and theme.

Tables present a comparison (either at 2100 or as average of 2040-2100) between a reference scenario and a comparing one. If the latter one is a random scenario run multiple times some basic statistic is given.

On the other side charts provide a quick overview of the dynamic of key variables over time.

This document is generated automatically trough a script directly from the model's output. Please refer to the individual articles for model description, scenario definition and output analysis.

\end{abstract}

\thispagestyle{empty}
\newpage
\tableofcontents
\thispagestyle{empty}
\newpage
\part[Tables]{Output tables}
\input{data/output/tables/00_master_tables}
\clearpage
\newpage
\part[Charts]{Output charts}

\clearpage
\subsection{Carbon charts}
\begin{figure}[htbp]
  \centering
  \caption{}
  \includegraphics[width=0.8\textwidth]{"charts/cbalance_ecoeffect_rpc00_FRA"}
  \label{fig:cbalance_ecoeffect_rpc00}
\end{figure}
\begin{figure}[htbp]
  \centering
  \caption{}
  \includegraphics[width=0.8\textwidth]{"charts/cbalance_ecoeffect_rpc86_FRA"}
  \label{fig:cbalance_ecoeffect_rpc86}
\end{figure}
\begin{figure}[htbp]
  \centering
  \caption{}
  \includegraphics[width=0.8\textwidth]{"charts/cbalance_cceffect_ext_FRA"}
  \label{fig:cbalance_cceffect_ext}
\end{figure}
\begin{figure}[htbp]
  \centering
  \caption{}
  \includegraphics[width=0.8\textwidth]{"charts/cbalance_cceffect_adapt_FRA"}
  \label{fig:cbalance_cceffect_adapt}
\end{figure}



%\newpage
%\bibliography{biblio}

\end{document}
\endinput

